\documentclass[11pt]{article}
\usepackage[margin=1in]{geometry}
\usepackage{times}
\usepackage{hyperref}
\hypersetup{colorlinks=true, linkcolor=black, urlcolor=blue, citecolor=black}
\title{High Voltage, Low Content: Lifters, Asymmetric Capacitors, and the Ion Wind Problem Revisited}
\author{Flyxion \\ \small Independent Researcher}
\date{}
\begin{document}
\maketitle

\begin{abstract}
Lifters, lightweight asymmetric-capacitor devices that rise off a benchtop when energized with high voltage, are a frequent hobbyist demonstration and an occasional exhibit in claims of electrogravitic propulsion. This paper argues that the phenomenon is explained by corona-discharge ionic wind, the same mechanism responsible for the Biefeld-Brown effect, and that the lifter's popularity as an antigravity demonstration owes to its visual resemblance to levitation rather than to any ambiguity in its underlying physics. The paper concentrates on what distinguishes the lifter from the larger capacitor literature: the engineering of the device, the scaling of its thrust with air density and power, and the use of the same mechanism in aircraft with no moving parts, which gives the ion wind account a test that a gravitational account could not pass. It also states the one point on which the older literature was uncertain, the apparent shortfall of simple momentum estimates, and reports how that point was settled.
\end{abstract}

\section{Introduction}

A lifter typically consists of a thin wire acting as the positive electrode, held a few centimeters above a wider foil skirt acting as the negative electrode, connected to a high-voltage, low-current power supply. When energized to several tens of kilovolts, the assembly rises and hovers, a demonstration popularized in hobbyist electronics communities and periodically resurfacing as evidence for reactionless or gravity-defying propulsion, generally accompanied by video of the device floating with no visible propeller or exhaust. The device is an asymmetrical capacitor thruster in the terminology of the technical reports, and the same configuration, thin wire above wide collector, is the standard geometry of the century-old Biefeld-Brown literature \cite{CMW2004}. A companion paper in this series treats the history of that effect and the arguments about its mechanism; the present paper takes the lifter as an engineered object and asks what it can and cannot do.

\section{Why This One Is Not Ambiguous}

Unlike some of the claims surveyed elsewhere in this series, the lifter is not a case of contested measurement or disputed replication. The mechanism is the production of ions at the thin wire by corona discharge, the acceleration of those ions toward the wider electrode, and the transfer of their momentum to the neutral air through repeated collisions. A paper in the Brazilian physics education literature describes the lifter as an asymmetric capacitor, explains the lift as the result of multiple collisions between ions and air molecules that transfer momentum, and estimates the lifting force from a simple model of the corona gap \cite{L2015}. The more rigorous analysis of Ianconescu and colleagues computes the levitation force from the unicharge ion wind and reports experimental results that confirm the calculation \cite{ISM2011}.

The device does not work in a vacuum. Brown's early reports of force in vacuum have not been reproduced under controlled conditions, and a bell jar test, which any hobbyist can attempt, leaves the lifter inert once the air is removed. That statement should be held to its proper strength: a failed hobbyist test is an anecdote, and the stronger evidence is the formal vacuum testing and the dependence of thrust on pressure and humidity discussed in the Biefeld-Brown literature \cite{T2004}. The conclusion is the same in both cases. Take away the air and the lifter stops, which is what a reaction against ionized air predicts and what a gravitational coupling does not.

\section{Scaling Gives the Decisive Test}

The most informative evidence on the lifter does not come from lifters at all. It comes from the engineers who took the same physics seriously as a propulsion method and asked how it scales. Gilmore and Barrett measured the performance of electroaerodynamic thrusters as a function of altitude and flight speed and showed that thrust and efficiency depend on the density of the surrounding air in the way ion-neutral momentum transfer requires \cite{GB2018}. A thruster that derived its force from a coupling to gravity would have no reason to weaken as the air thinned, so the altitude dependence is a direct discriminating observation, and the observation favors the ion wind account without reservation.

The same physics has since been made to carry load. Monrolin, Plouraboue and Praud developed the theory of electrohydrodynamic thrust for in-atmosphere propulsion, with an explicit model of how thrust scales with applied voltage, electrode spacing and current \cite{MPP2017}. Xu and colleagues then flew a fixed-wing aircraft with no moving parts in a gymnasium, using the thrust of an ionic wind from thin emitter wires toward downstream collector airfoils \cite{Xu2018}. The aircraft is a lifter-like structure that has been scaled and optimized, and it flies because the air it ionizes is pushed backward. The result matters here for a reason that has little to do with aircraft design. Antigravity claims commonly rest on the assertion that a lifter's thrust-to-power ratio is anomalous, exceeding what a simple wind model allows. The engineered thrusters offer a more useful comparison: their performance is described and optimized within the ion mobility and air density limits that the wind account sets, with no unexplained surplus invoked \cite{MPP2017,Xu2018}.

\section{Why the Antigravity Reading Persists Anyway}

The lifter's rhetorical power comes from what it looks like rather than what it does. A small, light, visually simple object rising silently off a table, with no propeller, no visible jet, and no moving parts, triggers the same intuitive surprise that made the Biefeld-Brown effect compelling a century earlier, and the surprise survives contact with the correct explanation because the correct explanation, ionized air pushing against more air, feels anticlimactic relative to the spectacle. This is a reliable pattern: demonstrations whose correct explanation is mundane but whose appearance is dramatic tend to retain their dramatic interpretation in public circulation long after the mundane explanation has been established in the specialist literature.

There is also a structural reason. The working fluid is invisible. A propeller or a rocket shows its reaction mass, and the observer therefore knows where the momentum goes. The lifter's reaction mass is the air column below and around it, moving slowly, in a low-velocity flow that the eye does not register. When an observer cannot see where the momentum goes, the absence of a visible exhaust is mistaken for the absence of an exhaust.

\section{The Entropic Gravity Framing}

On an account of gravitation as large-scale entropic and geometric bookkeeping rather than a locally sourced force, a lifter's electrostatic field has, as with the other devices surveyed in this series, no channel into gravitational curvature at all, since the field in question is an ordinary electric field acting on ionized air molecules, several orders removed in both scale and kind from anything that could plausibly redistribute the ordered configuration gravitation tracks. The lifter requires no appeal to this framework to be explained, which is itself informative: the mundane electrohydrodynamic account is sufficient on its own terms, and reaching for a gravitational reinterpretation adds an unnecessary and unsupported layer to an already complete explanation.

\section{Objections}

It is occasionally argued that thrust-to-power ratios for some lifter designs exceed what simple ion-drag models predict, suggesting an additional, uncharacterized contribution. The Army Research Laboratory report on the asymmetric capacitor did find that the ballistic ionic wind estimate fell short of the measured force by a large factor and concluded that the physical basis was not understood \cite{BF2003}. That discrepancy belongs to a particular crude estimate and not to the physics: later calculations that treat the ion drift and the resulting wind more carefully reproduce the measured force \cite{ISM2011}, and the electrohydrodynamic thrust theory developed for aircraft applications is built on the same ion wind picture \cite{MPP2017}. No lifter has been demonstrated to produce thrust in vacuum at any efficiency, which any genuinely novel non-electrohydrodynamic mechanism would presumably still permit.

A second objection is that the lifter works with either polarity of the supply, and that this is inconsistent with a simple ion wind. Polarity dependence is a standard feature of corona discharge, which differs between positive and negative emitters in onset voltage, current and ion species, and the Canning, Melcher and Winet study mapped the thrust of several asymmetrical capacitor thrusters as a function of voltage, polarity and grounding in exactly these terms \cite{CMW2004}. The dependence is part of the ion wind account and is not an anomaly requiring an additional one.

\section{Conclusion}

The lifter is among the more thoroughly explained devices in the antigravity literature. Its mechanism is ionic wind, its thrust depends on air density in the manner that mechanism requires, and the same physics has been scaled into aircraft that fly without moving parts. Its continued circulation as a gravity-defying demonstration says more about the gap between what an experiment looks like and what it measures than about any open question in the underlying physics. This paper presents no new measurements; it assembles published work and states that the one apparent anomaly in the early record has a documented resolution.

\begin{thebibliography}{99}

\bibitem{CMW2004} F. X. Canning, C. Melcher and E. Winet, Asymmetrical Capacitors for Propulsion, NASA Contractor Report NASA/CR-2004-213312, October 2004.

\bibitem{L2015} M. Cattani, A. Vannucci and V. G. Souza, Lifter: High Voltage Plasma Levitation Device, \emph{Revista Brasileira de Ensino de Fisica} 37 (2015) 3307, doi:10.1590/S1806-11173731869.

\bibitem{ISM2011} R. Ianconescu, D. Sohar and M. Mudrik, An Analysis of the Brown-Biefeld Effect, \emph{Journal of Electrostatics} 69 (2011) 512--521, arXiv:1011.1393.

\bibitem{T2004} M. Tajmar, Biefeld-Brown Effect: Misinterpretation of Corona Wind Phenomena, \emph{AIAA Journal} 42 (2004) 315--318.

\bibitem{GB2018} C. K. Gilmore and S. R. H. Barrett, Electroaerodynamic Thruster Performance as a Function of Altitude and Flight Speed, \emph{AIAA Journal} 56 (2018) 1105--1117.

\bibitem{MPP2017} N. Monrolin, F. Plouraboue and O. Praud, Electrohydrodynamic Thrust for In-Atmosphere Propulsion, \emph{AIAA Journal} 55 (2017) 4296--4305.

\bibitem{Xu2018} H. Xu, Y. He, K. L. Strobel, C. K. Gilmore, S. P. Kelley, C. C. Hennick, T. Sebastian, M. R. Woolston, D. J. Perreault and S. R. H. Barrett, Flight of an Aeroplane with Solid-State Propulsion, \emph{Nature} 563 (2018) 532--535.

\bibitem{BF2003} T. B. Bahder and C. Fazi, Force on an Asymmetric Capacitor, Army Research Laboratory Technical Report ARL-TR-3005, June 2003, arXiv:physics/0211001.

\end{thebibliography}

\end{document}
